% Introduction - Single author, direct style

Knowledge graphs underpin critical NLP applications: question answering systems retrieve facts from Freebase and Wikidata~\citep{bordes2015large}, dialogue agents ground responses in structured knowledge~\citep{moon2019opendialkg}, and fact verification systems check claims against knowledge bases~\citep{thorne2018fever}.
When these systems predict over unseen entities or novel relationships, knowing \emph{when to trust the output} is essential---a QA system should abstain rather than hallucinate, and a medical KG should flag uncertain drug interactions.
This motivates out-of-distribution (OOD) detection: distinguishing reliable predictions from uncertain ones.

Probabilistic knowledge graph embedding methods address this by learning distributions over entity embeddings rather than point estimates.
GP-KGE~\citep{Chen2021PERM} learns a variance $\sigma^2_e$ for each entity, capturing how ``uncertain'' the model is about that entity's representation.
The intuition is appealing: entities seen frequently should have low variance (well-constrained embeddings), while rare entities should have high variance.

\textbf{A fundamental limitation.}
This intuition breaks down because the learned variance is \emph{relation-agnostic}.
Consider an entity $e$ that appears frequently with relation $r_1$ (e.g., \texttt{born\_in}) but never with relation $r_2$ (e.g., \texttt{works\_at}).
GP-KGE assigns a single variance $\sigma^2_e$ to entity $e$, which is the same for both relations.
But the model's uncertainty about $(e, r_1, ?)$ and $(e, r_2, ?)$ should differ---$e$ has been observed with $r_1$ but not with $r_2$.

This observation leads to the central insight: \emph{uncertainty in knowledge graphs decomposes into semantic and structural components, and neither alone suffices for OOD detection.}
\textbf{Semantic uncertainty} (GP variance) captures embedding quality---rare entities have high variance.
\textbf{Structural uncertainty} (coverage) captures whether a specific entity-relation pair has been observed: $c(e, r) = 1$ if entity $e$ appears with relation $r$ in training.
Coverage is relation-specific by construction, addressing GP's limitation.

\textbf{Contributions.}
\begin{enumerate}
    \item A fundamental limitation in probabilistic KGE: learned variances are relation-agnostic and miss structural uncertainty.
    \item A semantic-structural decomposition framework that separates these orthogonal uncertainty types.
    \item A closed-form AUROC bound for coverage-based detection, validated empirically.
    \item Empirical demonstration that combining both signals (CAGP) achieves 0.87--0.96 AUROC across three benchmarks, with complementary failure modes validated on temporal and adversarial OOD settings.
\end{enumerate}

The key finding is not the CAGP algorithm itself---which is a simple linear combination---but the \emph{insight} that probabilistic KGE methods have a structural blind spot that cannot be learned away.
Figure~\ref{fig:main_results} summarizes the main results.

\begin{figure}[t]
\centering
\includegraphics[width=0.85\linewidth]{figures/fig1_main_results.pdf}
\vspace{-1em}
\caption{OOD detection AUROC. CAGP consistently outperforms single signals with 14--32\% synergy.}
\label{fig:main_results}
\end{figure}
